%% =====================================================================
%%  T-24-01  Refusing an Unauthorised Debt
%%  -------------------------------------------------------------------
%%  One idea per file. Core is mandatory. Slots in canonical order.
%%  Max 700 words. No layout commands. No filler. New source material
%%  for this entry goes in sources/<BOOK>/T-24-01.tex -- never by
%%  rewriting this file (docs/RULES.md R0.1).
%% =====================================================================
%% @kind: topic
%% @id: T-24-01
%% @title: Refusing an Unauthorised Debt
%% @subtitle: Declining the premise rather than arguing the amount
%% @chapter: c24
%% @order: 10
%% @status: stable
%% @sources: B4
%% @updated: 2026-09-19
%% @allow-rewrite: slot migration to the seven-question format (docs/RULES.md section 2)

\topic{T-24-01}{Refusing an Unauthorised Debt}{Declining the premise rather than arguing the amount}


\begin{core}
The answer to an obligation you did not create is not to argue that you owe nothing. It is to decline the premise on which the obligation was built --- promptly, politely, and without engaging the merits of the request that follows it.
\end{core}

\begin{mechanism}
Arguing about the size of a debt concedes that a debt exists. The obligation is installed by the occurrence of a benefit rather than by its value, so any discussion of value takes place inside the frame the other party built and can only ever adjust the amount.

What removes the frame is a decision taken before the request arrives: that an unrequested benefit opens no account. That decision is easy to hold and hard to execute, because the pressure is social and visible, and refusing looks like ingratitude to any audience. The refusal therefore has two halves --- an internal reclassification, and an external form of words that does not turn the refusal into a verdict on the giver's character. Where the benefit genuinely cannot be refused, only the second half changes: repay it immediately, specifically, and in a currency of your own choosing, so that nothing remains outstanding to be drawn on later.
\end{mechanism}

\begin{application}
  \item Decide the policy in advance, so that the decision is not made under pressure in front of an audience.
  \item Decline in the same breath as the offer, before acceptance has been seen by anyone.
  \item Where decline is impossible, repay at once, name the amount, and state that the account is closed.
  \item Answer the request separately and explicitly: these are two questions, and the first does not settle the second.
  \item Keep the refusal short, warm and unexplained.
\end{application}

\begin{feedback}
  \item A benefit you did not request has arrived, and the giver chose its form and timing.
  \item A request follows that has no connection to the benefit.
  \item You are aware of feeling obliged before you have decided anything.
  \item Refusing would be observed by a third party.
  \item The benefit is referred to in the same conversation as the request.
\end{feedback}

\begin{failure}
A refusal that is habitually applied makes you difficult to help, and people stop offering --- which removes the ordinary goodwill that most relationships run on. The policy is therefore situational rather than absolute: it is aimed at benefits that were unrequested, chosen by the giver, and followed quickly by an unrelated ask. Applied to genuine kindness it produces a person who is defended against nothing and isolated by the defence.
\end{failure}

\begin{countermeasures}
  \item When the refusal is called ungrateful, agree that you are grateful and decline anyway. Gratitude and compliance are different things and can be stated separately.
  \item When the benefit is re-described as something you wanted, do not dispute the description; restate that you did not ask for it.
  \item When a third party applies the pressure, address them rather than the giver: the account, if there is one, is not theirs to collect.
  \item When the request is reduced after your refusal, treat the reduction as a new proposal and evaluate it on its own --- see T-16-01.
  \item When you are offered a way to repay, take it, pay it in full, and consider the matter finished rather than continuing to negotiate.
\end{countermeasures}

\sources{B4}
\seesources
\seealso{T-19-01, T-19-02, T-16-01, T-04-02}
